\subsection{A stopped hedge portion leaves a reserve behind}
\label{sec:partial-hedge-discontinuation}

Return to the original forecast purchase of 100 tonnes at an initial forward
price of 80. At the first observation, price is 90 and the derivative has
gained 1200: the cash flow hedge reserve is 1000 and cumulative
ineffectiveness is 200. Now management finds that 20 tonnes are no longer
highly probable, while the remaining 80 tonnes continue to satisfy the hedge
conditions. The derivative is still owned. Stopping hedge accounting for a
portion therefore changes where its later gains go, even if no contract is
sold and no cash changes hands.

IFRS 9.B6.5.27(b) requires partial discontinuation for the forecast volume
that is no longer highly probable; paragraphs 6.5.6 and B6.5.22 make the
change prospective. Paragraph B6.5.23 prevents voluntary discontinuation of
a relationship that continues to qualify. These are the inspected historical
EU 2023/1803 provisions \citep{eu20231803}. The example supplies the
qualification and documentation findings; a percentage reduction alone cannot
establish them.

For this case, the 100 tonnes have identical per-unit exposures, one original
measurement date and one price origin. The identified 20 tonnes correspond
to 20\% of the entire derivative, whose proportional designation is permitted
by paragraph 6.2.4(c). This homogeneity matters. Multiplying both sides of the
lower-of calculation by a nonnegative fraction \(a\) gives
\[
 \min\{a|D|,a|H|\}=a\min\{|D|,|H|\},
\]
where \(D\) is cumulative derivative gain and \(H\) the cumulative change
in the hedged cash flows. With \(a=0.2\), the existing amounts divide as
follows, without creating another gain or resetting the original price:

\begin{center}
\small
\begin{tabular}{lrrr}
\toprule
& Continuing 80 tonnes & Stopped 20 tonnes & Total\\
\midrule
Derivative value and gain & 960 & 240 & 1200\\
Cash flow hedge reserve & 800 & 200 & 1000\\
Prior ineffectiveness & 160 & 40 & 200\\
\bottomrule
\end{tabular}
\end{center}

The stopped portion now carries two distinct origins: derivative value 240
and accumulated reserve 200. Suppose a later observation gives price 95
and whole-derivative value 1400. The continuing derivative portion has gained
\(0.8\times1400=1120\) since original designation; its hedged cash-flow
change is \(-80(95-80)=-1200\). Its reserve is therefore 1120 and its
cumulative ineffectiveness is zero. The earlier continuing ineffectiveness
of 160 has reversed.

The stopped derivative portion is now worth \(0.2\times1400=280\).
Its additional gain of \(280-240=40\) enters profit under paragraph 5.7.1;
it does not increase the stopped reserve. Alongside the prior stopped
ineffectiveness of 40, it leaves total cumulative profit of 80. If the
stopped purchase is still expected, total reserve is
\(1120+200=1320\). Applying the lower-of rule to the whole original
position would instead give reserve 1400 and profit zero. That calculation
would incorrectly keep hedging the portion whose designation has ended.

The distinction between ``highly probable'' and ``expected'' decides what
happens to the frozen 200. Paragraph 6.5.12 explicitly allows a forecast
transaction to remain expected after it ceases to be highly probable.
Its reserve then remains until the cash flows occur, subject to the
unrecoverable-loss rule. If the stopped purchase is no longer expected,
the remaining 200 moves immediately from the reserve to profit. Total
reserve and profit become 1120 and 280. A query or later valuation does
not itself make that expectation judgment.

Follow both outcomes into inventory, using the supplied later price and
settlement values. If all 100 tonnes arrive, supplier payment is 9500
and the remaining reserve of 1320 reduces initial inventory cost to 8180.
With opening cash and capital of 20000, receiving derivative settlement
of 1400 leaves cash \(20000-9500+1400=11900\). Cash plus inventory
is 20080, equal to capital plus profit 80. If the stopped 20 tonnes are
canceled, only 80 tonnes arrive: payment 7600 less reserve 1120 gives
inventory 6480, while cash is 13800. Net assets of 20280 then equal
capital plus profit 280. The released reserve cannot also reduce inventory.
Supplier and derivative cash movements remain separately identified.

A loss reserve needs an additional check. Reverse the price path: the first
observation is 70 with derivative value \(-1200\), followed by price 65
and value \(-1400\). Before any loss-recovery adjustment, continuing and
stopped reserves are \(-1120\) and \(-200\), with total profit
\(-80\). Suppose new evidence shows that 60 of the stopped loss will
not be recovered. Paragraph 6.5.11(d)(iii) requires that 60 to enter
profit immediately, leaving stopped reserve \(-140\) and total profit
\(-140\). If the full purchase then occurs, payment 6500 plus the
remaining total loss reserve of 1260 gives inventory cost 7760. If the
stopped purchase is canceled instead, only its remaining loss of 140 is
released. Charging the original 200 again would count the earlier 60 twice.

This allocation follows from identified homogeneous portions. Different
price origins, unmatched instrument reductions or a heterogeneous reserve
need their own allocation analysis. The example also assumes an eligible
external non-option derivative, supplied complete valuations, recoverability
before the specified loss finding, no fees and a pretax reporting scope.
Those facts make the reserve history reconstructible; they do not provide
a pricing model or a general rule for splitting any accumulated reserve.
